\section{问题二的分析与求解：减重维持失败因素识别}

\subsection{研究目标与建模边界}

问题二要求识别减肥后体重反弹或治疗中断的主要原因。由于现有证据同时包含行为因素、心理因素、支持环境和生理代偿，若直接将不同结局、不同统计口径的数据混合排序，会使结果失去可比性。为此，本文将研究分为两个层次：第一层从系统综述中筛选15项具有统一证据口径的减重维持因素，采用CRITIC客观赋权与TOPSIS逼近理想解排序；第二层利用依从性、治疗退出、生理代偿和中国肥胖管理调查等资料解释减肥失败的作用机制。

正式排序仅反映文献证据的综合优先级，不表示某一因素对个体减肥失败的因果概率。时间不足、家庭支持、治疗退出、疾病、生理代偿及ACTION-China调查中的管理障碍均只用于机制分析，不进入15项因素的TOPSIS排序。

\subsection{证据来源与纳入规则}

\subsubsection{证据来源}

15项正式排序因素来自Varkevisser等对67篇研究、124项减重维持决定因素的系统综述\textsuperscript{[10]}。Lemstra等的依从性Meta分析\textsuperscript{[11]}、治疗退出综述\textsuperscript{[12]}、生理代偿综述\textsuperscript{[13]}、ACTION-China调查\textsuperscript{[14]}及国家指南\textsuperscript{[15--16]}只用于解释长期坚持、退出、医学反应和中国管理情境，不进入排序。

\subsubsection{因素筛选与指标统一}

正式排序需同时满足三个条件：一是属于减重后可干预的行为、监测或心理因素；二是原综述给出了研究数量、方向判断和证据等级；三是能够统一改写为“缺少该因素会提高维持失败风险”。年龄、性别等不可干预特征，以及方向不一致的一般社会支持不进入排序。疾病、妊娠、身体限制和生理代偿虽具有现实意义，但缺少与15项因素同口径的三指标数据，因此仅在机制部分讨论。

设第$i$项因素在第$j$个评价指标上的原始值为$x_{ij}$。三个效益型指标分别为研究数量$n_i$、统一方向后的一致率和证据等级，其中Strong编码为2，Moderate编码为1。完整15项决策矩阵保存在\path{q2_factor_evidence.csv}，正文只展示最终排序。

需要特别说明的是，M4来自7篇研究，而非8篇研究。原综述将同一研究中的不同不健康食品指标分别判断，共形成8次方向判断，其中7次支持、1次无效，因此一致率为$7/8=0.875$；模型中的研究数量仍取$n_4=7$。

\subsection{CRITIC客观赋权--TOPSIS排序模型}

\subsubsection{CRITIC客观赋权}

CRITIC方法根据指标的变异程度和指标间冲突程度确定客观权重\textsuperscript{[17]}。由于三个指标均为效益型，先进行极差标准化：
\begin{equation}
z_{ij}^{(C)}=
\frac{x_{ij}-\min_i x_{ij}}
{\max_i x_{ij}-\min_i x_{ij}} .
\end{equation}
式中，$z_{ij}^{(C)}$为第$i$项因素在第$j$个指标上的CRITIC标准化值，$x_{ij}$为原始指标值。

第$j$个指标的离散程度用标准差表示：
\begin{equation}
\sigma_j=
\sqrt{\frac{1}{m}\sum_{i=1}^{m}
\left(z_{ij}^{(C)}-\bar z_j^{(C)}\right)^2}.
\end{equation}
式中，$m=15$为因素数量，$\bar z_j^{(C)}$为第$j$个指标标准化值的均值，$\sigma_j$越大表示该指标对不同因素的区分能力越强。

指标$j$与指标$k$之间的Pearson相关系数为
\begin{equation}
r_{jk}=
\frac{\sum_{i=1}^{m}
\left(z_{ij}^{(C)}-\bar z_j^{(C)}\right)
\left(z_{ik}^{(C)}-\bar z_k^{(C)}\right)}
{\sqrt{\sum_{i=1}^{m}\left(z_{ij}^{(C)}-\bar z_j^{(C)}\right)^2}
\sqrt{\sum_{i=1}^{m}\left(z_{ik}^{(C)}-\bar z_k^{(C)}\right)^2}} .
\end{equation}
式中，$r_{jk}$越小，说明两个指标提供的重复信息越少，冲突性越强。

综合离散程度和冲突程度，第$j$个指标的信息量及权重分别为
\begin{equation}
C_j=\sigma_j\sum_{k=1}^{p}(1-r_{jk}),
\end{equation}
\begin{equation}
w_j=\frac{C_j}{\sum_{j=1}^{p}C_j}.
\end{equation}
式中，$p=3$为指标数量，$C_j$为信息量，$w_j$为CRITIC客观权重，且$\sum_{j=1}^{p}w_j=1$。

\subsubsection{TOPSIS逼近理想解排序}

TOPSIS通过比较各因素与正、负理想解的距离进行综合排序\textsuperscript{[18]}。首先对原始决策矩阵进行向量标准化：
\begin{equation}
z_{ij}^{(T)}=
\frac{x_{ij}}{\sqrt{\sum_{i=1}^{m}x_{ij}^{2}}}.
\end{equation}
式中，$z_{ij}^{(T)}$为TOPSIS向量标准化值。

将CRITIC权重引入标准化矩阵，得到加权决策矩阵：
\begin{equation}
v_{ij}=w_j z_{ij}^{(T)}.
\end{equation}
式中，$v_{ij}$为第$i$项因素在第$j$个指标上的加权标准化值。

由于三个指标均为效益型，正、负理想解分别为
\begin{equation}
A^+=\left\{\max_i v_{ij}\mid j=1,\ldots,p\right\},
\qquad
A^-=\left\{\min_i v_{ij}\mid j=1,\ldots,p\right\}.
\end{equation}
式中，$A^+$表示各指标最优值组成的理想方案，$A^-$表示各指标最差值组成的负理想方案。

第$i$项因素到正、负理想解的欧氏距离为
\begin{equation}
D_i^+=\sqrt{\sum_{j=1}^{p}\left(v_{ij}-A_j^+\right)^2},
\qquad
D_i^-=\sqrt{\sum_{j=1}^{p}\left(v_{ij}-A_j^-\right)^2}.
\end{equation}
式中，$D_i^+$和$D_i^-$分别表示因素$i$与最好、最差证据组合之间的距离。

最终贴近度为
\begin{equation}
S_i=\frac{D_i^-}{D_i^++D_i^-},\qquad 0\leq S_i\leq 1.
\end{equation}
式中，$S_i$越大，表示该因素的综合证据越接近理想水平。对$S_i$降序排列，并对相同得分采用$1,2,2,4$形式的竞赛排名。

\subsection{基准结果与证据优先级}

CRITIC计算得到研究数量、方向一致率和证据等级的权重依次为
\[
0.238194473347,\qquad 0.419561680818,\qquad 0.342243845835.
\]
方向一致率权重最高，说明在当前15项因素中，该指标同时具有较强区分度和较多非重复信息。各因素的TOPSIS得分与竞赛排名如表~\ref{tab:q2-ranking}~所示。

\begin{table}[H]
\centering
\small
\caption{CRITIC权重与15项因素TOPSIS排序结果}
\label{tab:q2-ranking}
\setlength{\tabcolsep}{4pt}
\begin{tabular}{ccp{7.4cm}cc}
\toprule
名次 & 编号 & 因素 & TOPSIS得分 & 类别 \\
\midrule
1 & M1 & 体力活动没有持续增加 & 0.838179 & 运动 \\
2 & M2 & 未持续降低总能量摄入 & 0.466327 & 饮食 \\
2 & M3 & 不经常称重或监测体重 & 0.466327 & 自我监测 \\
4 & M4 & 未减少甜食、油炸、快餐等不健康食品 & 0.362655 & 饮食 \\
5 & M5 & 果蔬摄入没有增加 & 0.324335 & 饮食 \\
6 & M6 & 含糖饮料没有减少 & 0.292081 & 饮食 \\
6 & M7 & 运动自我效能低或感到运动障碍大 & 0.292081 & 心理 \\
8 & M13 & 体重管理自我效能低 & 0.287990 & 心理 \\
9 & M15 & 内在线索引起的失控进食严重 & 0.275129 & 心理 \\
10 & M9 & 没有控制食物份量 & 0.265039 & 饮食 \\
10 & M10 & 没有持续记录饮食 & 0.265039 & 自我监测 \\
10 & M11 & 没有降低膳食脂肪摄入 & 0.265039 & 饮食 \\
13 & M14 & 饮食自我效能低 & 0.168554 & 心理 \\
14 & M8 & 暴食、情绪性或失控性进食 & 0.151633 & 心理 \\
15 & M12 & 没有监测运动 & 0.097721 & 自我监测 \\
\midrule
\multicolumn{3}{r}{CRITIC权重：研究数量/一致率/证据等级} &
\multicolumn{2}{l}{0.238194/0.419562/0.342244} \\
\bottomrule
\end{tabular}
\end{table}

基准方案的前五项为：M1排名第一，M2与M3并列第二，M4排名第四，M5排名第五。M1得分为0.838179，明显高于并列第二的0.466327，表明“体力活动没有持续增加”在研究数量、方向一致率和证据等级的综合意义下最受现有证据支持。饮食热量控制与体重监测并列第二，说明减重维持需要同时关注能量摄入和持续反馈，而不能仅依赖单次减重结果。

表~\ref{tab:q2-ranking}~给出精确分值，图~\ref{fig:q2-ranking}~展示类别与名次；二者均表示证据优先级，不可解释为个体发生减肥失败的概率。

\begin{figure}[H]
\centering
\includegraphics[width=0.92\textwidth]{减重维持因素排行.png}
\caption{15项减重维持因素的CRITIC--TOPSIS排序}
\label{fig:q2-ranking}
\end{figure}

\subsection{减肥失败的四层作用机制}

排序结果说明了可比因素的证据优先级，但减肥失败并非单一行为造成。结合补充证据，可将其概括为行为执行、心理调节、支持环境和生理医学四个相互作用的层次。

\textbf{（1）行为执行层。} M1、M2、M4和M5影响长期能量缺口，M3提供反馈纠偏。Meta分析中总体依从率为60.5\%，干预不足12个月时为69.9\%，至少12个月时降至53.0\%\textsuperscript{[11]}，说明维持期延长伴随更多行为中断。

\textbf{（2）心理调节层。} 自我效能和情绪性进食会影响计划执行；治疗退出综述报告的退出率为5\%--62\%，主要原因涉及个人问题、治疗不满意、健康状况、动机、物流和时间\textsuperscript{[12]}，不能把失败简单归因于意志力。

\textbf{（3）支持环境层。} 监督式与自我监测式项目依从率分别为68.6\%和41.5\%，社会支持使依从率比提高至1.29\textsuperscript{[11]}。ACTION-China中仅40.8\%的肥胖者近5年与专业人员讨论过体重，34.8\%曾获专业诊断\textsuperscript{[14]}，说明监督、支持和医患沟通会影响长期执行。

\textbf{（4）生理医学层。} 每减重1\,kg，每日能量消耗可额外下降20--30\,kcal，食欲驱动约增加100\,kcal/日；长期研究中5年内反弹超过80\%\textsuperscript{[13]}。若生活方式干预3--6个月仍不能减重5\%，指南建议进一步评估并调整治疗\textsuperscript{[15--16]}。

时间、家庭支持、治疗退出、疾病、生理代偿和ACTION-China管理障碍的结局或统计口径与15项因素不同，故均不进入TOPSIS计算，只用于解释减肥失败如何形成和放大。

\subsection{六场景灵敏度分析}

为检验排序对一致率估计、研究数量变换和权重设定的敏感程度，构造六种场景：基准CRITIC、Wilson 95\%下限一致率、Laplace修正一致率、研究数取$\log(1+n)$、TOPSIS等权重和删除证据等级。Spearman等级相关系数采用并列值的平均秩计算，论文展示名次仍采用竞赛排名。结果如表~\ref{tab:q2-sensitivity}~所示。

\begin{table}[H]
\centering
\scriptsize
\caption{六种场景下的排序灵敏度分析}
\label{tab:q2-sensitivity}
\resizebox{\textwidth}{!}{
\begin{tabular}{p{4.0cm}cccccc}
\toprule
场景 & 研究数权重 & 一致率权重 & 证据权重 & Spearman系数 & 前五重合 & 最大名次变化 \\
\midrule
原始一致率+CRITIC & 0.238194 & 0.419562 & 0.342244 & 1.000000 & 5 & 0 \\
Wilson 95\%下限 & 0.216450 & 0.280130 & 0.503420 & 0.967509 & 5 & 3 \\
Laplace修正 & 0.245670 & 0.348971 & 0.405359 & 0.935018 & 5 & 3 \\
研究数取$\log(1+n)$ & 0.249887 & 0.429099 & 0.321014 & 1.000000 & 5 & 0 \\
TOPSIS等权重 & 0.333333 & 0.333333 & 0.333333 & 0.974729 & 5 & 3 \\
删除证据等级 & 0.400316 & 0.599684 & 不适用 & 0.865245 & 5 & 7 \\
\bottomrule
\end{tabular}}
\end{table}

六种场景的前五项均为M1--M5，前五重合数量均为5。Spearman系数满足
\[
0.865245356821\leq \rho_s\leq 1.000000000000,
\]
说明核心排序具有较强稳健性。删除证据等级时最大名次变化达到7，主要影响中后部因素，表明这些因素对Strong/Moderate编码较敏感。该波动不改变前五结论，但提示中后部得分接近时不宜作过度精细的先后解释。

\subsection{模型局限性与本题结论}

\textbf{模型局限性。} 主综述没有提供可统一合并的效应量，研究数量和方向一致率不能替代样本量与效应大小；证据等级又部分利用研究数量和一致性，可能造成重复计权。因此本文只解释前五项及其稳健性，不对接近的中后部名次作精细因果解释。ACTION-China和生理代偿证据因口径不同，仅用于机制说明。

\textbf{本题结论。} 基于15项同口径因素，CRITIC赋予研究数量、一致率和证据等级的权重分别为$0.238194473347$、$0.419561680818$和$0.342243845835$。TOPSIS结果显示，M1“体力活动没有持续增加”排名第一；M2“未持续降低总能量摄入”和M3“不经常称重或监测体重”并列第二；M4排名第四，M5排名第五。六种灵敏度场景的前五项完全重合，最低Spearman系数为$0.865245356821$。因此，减重维持应优先保障持续运动、能量摄入控制、体重监测及饮食结构调整，同时通过心理支持、专业监督和医学评估降低治疗退出与长期反弹风险。
